% revizyon_kontrolleri.tex — PRAB revizyonu öncesi geçerlilik kontrolleri (Adım 0–6, T1, T2)
% Derleme: xelatex revizyon_kontrolleri.tex (iki kez)
\documentclass[11pt,a4paper]{article}
\usepackage{fontspec}
\usepackage{polyglossia}
\setdefaultlanguage{turkish}
\setotherlanguage{english}
\usepackage{amsmath,amssymb}
\usepackage{graphicx}
\usepackage{booktabs}
\usepackage{array}
\usepackage{tabularx}
\usepackage[margin=2.3cm]{geometry}
\usepackage{xcolor}
\usepackage{hyperref}
\usepackage{caption}
\usepackage{float}
\hypersetup{colorlinks=true,linkcolor=blue!50!black,urlcolor=blue!50!black,citecolor=blue!50!black}
\captionsetup{font=small,labelfont=bf}
\graphicspath{{sonuclar/}}
\setlength{\parskip}{0.45em}
\setlength{\parindent}{0pt}

% Türkçe dil ayarı \texttt içinde "=" çevresine boşluk ekler; kod makrosunda kapat
\newcommand{\kod}[1]{{\XeTeXinterchartokenstate=0\texttt{#1}}}
\newcommand{\BEL}{\textsc{[belgelenmiş]}}
\newcommand{\YOR}{\textsc{[yorum]}}
\newcommand{\HIP}{\textsc{[hipotez]}}
\newcommand{\Csym}{C}

\title{Ters Yönde Dönen Kapalı Yörüngelerle Simetrik Kaçıklık İzleme:\\
PRAB Revizyonu Öncesi Geçerlilik Kontrolleri}
\author{Hazırlayan: Claude (Anthropic yapay zekâ asistanı)\\
\small arXiv:2608.04647v1 ve hakem raporları üzerine yapılan simülasyon kontrolleri}
\date{25 Ekim 2026}

\begin{document}
\maketitle

\begin{abstract}
\noindent Bu belge, \emph{``Accessing the symmetric component of quadrupole misalignment through
counter-rotating closed orbits in EDM storage rings''} makalesine gelen iki hakem raporundan sonra yapılan
yedi grup simülasyon kontrolünü anlatır. Amaç makaleyi savunmak değil, yöntemin \emph{nerede geçerli
olduğunu} bulmaktır. Kısa sonuç şudur. Makalede yazıldığı haliyle yöntem (BPM'ler laboratuvara sabit
varsayılır, fit yalnızca kuadrupol sütunları içerir), BPM'ler yer hareketiyle birlikte kaydığında ve
saptırıcı (deflector) alanları sürüklendiğinde çalışmaz. Buna karşılık üç tasarım koşulu altında yöntem
güçlü bir biçimde ayakta kalır: (i)~BPM'ler kuadrupollere rijit bağlıdır ve kestirici bunu bilir,
(ii)~düzelticiler kuadrupol üstüne sarılmış manyetik trim sargılarıdır, (iii)~saptırıcılar kendi elektrik
trimleriyle düzeltilir. Bu koşullarda iki demetli izleme tek demetten 25--50 kat iyi sonuç verir ve bir
hafta boyunca sahte EDM'yi BBA (demet tabanlı hizalama) kalitesinde tutar. Ayrıca makaledeki ``simetrik
bileşen yörüngeye kördür'' iddiasının yalnızca düşük harmoniklerde ve 4--7 kat düzeyinde geçerli olduğu
bulunmuştur.
\end{abstract}

\tableofcontents
\newpage

%======================================================================
\section{Bu belge nasıl okunmalı}
%======================================================================

Bu bölümde belgenin yapısını, kullanılan büyüklükleri ve iddiaların nasıl etiketlendiğini anlatıyorum.

\subsection{İddia etiketleri}
Belgedeki her önemli iddia üç türden biriyle işaretlenmiştir:
\begin{itemize}
\item \BEL: Bu çalışmada gerçekten çalıştırılmış bir simülasyonun çıktısı ya da bir kaynaktan doğrudan
alınmış bilgi. Kaynağı (betik ve sonuç dosyası) yanında verilir.
\item \YOR: Belgelenmiş sonuçlardan benim çıkardığım sonuç. Mantığı açıkça yazılır.
\item \HIP: Test edilmemiş bir tahmin ya da varsayılmış bir parametre değeri.
\end{itemize}

\subsection{Temel büyüklükler}

\paragraph{Simetrik ve antisimetrik bileşen.} Halkada 24 FODO hücresi, yani 48 kuadrupol vardır: her
hücrede bir odaklayıcı (QF) ve bir odaklamayan (QD). Hücre $k$'daki iki kuadrupolün kaçıklıkları
$v_{2k}$ (QF) ve $v_{2k+1}$ (QD) olsun. Bu iki sayı her zaman bir ortalamaya ve bir farka ayrılabilir:
\begin{equation}
v^{s}_{2k}=v^{s}_{2k+1}=\tfrac12\,(v_{2k}+v_{2k+1}),\qquad
v^{a}_{2k}=-v^{a}_{2k+1}=\tfrac12\,(v_{2k}-v_{2k+1}).
\end{equation}
Burada $v^s$ \emph{simetrik} bileşendir: iki kuadrupol aynı yöne kaymıştır. $v^a$ \emph{antisimetrik}
bileşendir: iki kuadrupol zıt yönlere kaymıştır. Her kaçıklık deseni bu ikisinin toplamıdır,
$v=v^s+v^a$. Birimleri metredir; tablolarda mikrometre ($\mu$m) kullanılır.

\paragraph{Sahte EDM ve hedef.} Sahte EDM, spinin dikey bileşeninin büyüme hızıdır, $f=\mathrm{d}S_y/\mathrm{d}t$
(birimi rad/s). Deneyin hedefi $10^{-29}\,e\cdot$cm'lik bir EDM'dir; bu simülasyonda
$f\approx 1$~nrad/s'ye karşılık gelir. Belgedeki bütün sahte EDM değerleri bu hedefin katı olarak verilir:
``$C=3$'' demek, hedefin üç katı demektir.

\paragraph{İki demetli metrik.} Gerçek EDM, demet yönü çevrilince (saat yönü CW'den saat yönünün tersi
CCW'ye) işaret değiştirir. Bu yüzden deneyde ölçülen büyüklük farktır:
\begin{equation}
C=\tfrac12\,\bigl|f_{\rm CW}-f_{\rm CCW}\bigr| .
\end{equation}
İki yönde aynı işaretli kalan kısım, yani $\tfrac12(f_{\rm CW}+f_{\rm CCW})$, bu farkta sadeleşir. Ona
``çift kısım'' diyorum. Makaledeki şekillerde kullanılan metrik de $C$'dir.

\paragraph{Tepki matrisi ve kestirici.} $R$, kuadrupol kaçıklığından BPM okumasına giden $48\times48$
matristir: $y=R\,v$. CW ve CCW için ayrı matrisler vardır ($R_{\rm CW}$, $R_{\rm CCW}$); C++ izleyiciyle
ölçülmüşlerdir. Kaçıklığı okumadan geri bulmak için \emph{kesik tekil değer ayrışımı} (TSVD, truncated
singular value decomposition) kullanılır: matrisin en büyük tekil değerinin \kod{rcond} katından küçük
tekil değerleri atılır, böylece gürültünün büyütülmesi önlenir.

\paragraph{Hesap araçları.} Bütün spin sonuçları repodaki GL4 semplektik parçacık ve Thomas--BMT spin
izleyicisiyle (\kod{integrator.cpp}) elde edilmiştir. Saptırıcı hataları ve saptırıcı konumundaki
düzelticiler için bu motorun değiştirilmemiş bir kopyası yazılmıştır (\kod{revizyon/integrator\_rev.cpp}).
Hatalar sıfırken yeni motor eskisiyle birebir aynı yörüngeyi verir (fark 0)~\BEL. Motorun doğruluğu,
kayıtlı bir sonucun yeniden üretilmesiyle de kontrol edildi: seed~0, $\sigma=10\,\mu$m için kayıtlı
$f=1.0890\times10^{-6}$~rad/s birebir tekrar elde edildi~\BEL.

%======================================================================
\section{Adım 0: Kodda ne varsayılmış?}
%======================================================================

Bu adımda hiçbir şey değiştirmeden mevcut kodun hangi fiziksel varsayımları yaptığını inceledim. Sonuçlar,
sonraki bütün adımların neden gerekli olduğunu açıklıyor.

\begin{enumerate}
\item \textbf{BPM'ler laboratuvara sabit.} BPM okuması ``yörünge + statik ofset + gürültü'' olarak
modellenmiş; kuadrupol konumu okumadan çıkarılmıyor (\kod{kmod\_drivers/twobeam\_oc.py}, satır 96--106)~\BEL.
Makalenin Sec.~V.B'deki ``BPM kuadrupolle birlikte hareket eder'' cümlesi hiçbir simülasyonda yok.
\item \textbf{Düzeltici, kuadrupol içindeki manyetik dipol.} Düzeltme, kestirilen kaçıklığın doğrudan
kuadrupol merkezinden çıkarılmasıyla yapılıyor. Bu, kuadrupolün tam içine yerleştirilmiş bir manyetik
düzelticiye eşdeğer~\BEL. Makale düzeltici türünü yazmıyor.
\item \textbf{Saptırıcı alan hatası yok; roll hatası yanlış fizikle var.} Mevcut \kod{dipole\_tilt}
parametresi elektrik alanını döndürmüyor, radyal bir \emph{manyetik} alan ekliyor
(\kod{integrator.cpp}:153--170)~\BEL. Manyetik kuvvet hız yönüne bağlı olduğu için bu model CCW demetinde
işaret değiştirir; gerçek elektrik roll'u ise değiştirmez.
\item \textbf{Hizalama hatası korelasyonsuz Gauss.} Hücre~0'daki QF'in kaçıklığı her zaman sıfır
alınmış~\BEL.
\item \textbf{İki demetli fit:} $S=[R_{\rm CCW};R_{\rm CW}]$, \kod{rcond}~$=0.01$ ile TSVD; makalede
``regularized pseudo-inverse'' diye geçiyor~\BEL.
\item \textbf{Şekil 4'teki elmas} (iki demetli düzeltme) 14~nm BPM gürültüsü ve \emph{sıfır} BPM ofsetiyle
hesaplanmış. 100~$\mu$m BPM ofsetiyle aynı hesap hedefin medyan 198 katını veriyor
(\kod{twobeam\_oc\_14nm\_g20\_boff100um.json})~\BEL. Şekil açıklaması 5 seed diyor, elmas ise 15 seed'den geliyor.
\end{enumerate}

%======================================================================
\section{Adım 1: BPM modeli}
%======================================================================

Bu bölümde makaledeki olası bir iç tutarsızlığı test ediyorum: Sec.~III'teki ``simetrik bileşen
yörüngeye kör'' argümanı laboratuvara sabit BPM varsayar; Sec.~V.B ise BPM'lerin kuadrupolle birlikte
hareket ettiğini söyler.

\subsection{Sezgi}
Bir kuadrupol $d$ kadar kaysın. Laboratuvara sabit bir BPM, demetin yörüngesini okur; simetrik bir kaymanın
yörüngede bıraktığı iz küçüktür. Ama BPM kuadrupole bağlıysa, kuadrupolle birlikte $d$ kadar kayar. Demet
yerinden pek oynamadığı için BPM bu sefer yaklaşık $-d$ okur. Yani simetrik kayma, tek demette bile
büyük bir sinyal olarak görünür. Bu durumda iki demete gerek kalmayabilir.

\subsection{Model}
BBA kalibrasyonundan sonraki okuma değişimi şöyle yazılır:
\begin{equation}
\Delta r_b = R_b\,\Delta m - u + n_b ,\qquad u=\rho\,\Delta m .
\end{equation}
Burada $\Delta m$ kuadrupol manyetik merkezlerinin sürüklenmesi (bilinmeyen), $u$ BPM'in elektriksel
merkezinin hareketi, $n_b$ demet $b$'nin BPM gürültüsü, $\rho$ ise BPM'in kuadrupolü ne kadar izlediğidir.
$\rho=0$ laboratuvara sabit BPM, $\rho=1$ rijit bağlı BPM demektir. Dört kestirici karşılaştırıldı:

\begin{center}\small
\begin{tabular}{lll}
\toprule
& demet sayısı & varsayım \\
\midrule
E1 & tek demet & BPM sabit (standart yörünge düzeltmesi) \\
E2 & tek demet & $\rho$ biliniyor: $(R-\rho I)$ \\
E3 & iki demet & BPM sabit (makaledeki yöntem) \\
E4 & iki demet & $\rho$ biliniyor: $[R_{\rm CCW}-\rho I;\,R_{\rm CW}-\rho I]$ \\
\bottomrule
\end{tabular}
\end{center}

Her kestiricinin TSVD eşiği bir ızgarada taranıp kendi en iyi değeriyle değerlendirildi; makaledeki sabit
\kod{rcond}~$=0.01$ de ayrıca raporlandı. Senaryo: 5~$\mu$m rms saf simetrik sürüklenme, 200 seed
(yeniden kurma) ve 3 seed (C++ sahte EDM).

\subsection{Sonuçlar}

\begin{table}[H]\centering\small
\caption{Kalan simetrik rms [$\mu$m], 14~nm gürültü, 200 seed medyanı. Başlangıç 5~$\mu$m.
\kod{revizyon/sonuclar/adim1\_recon.json}~\BEL.}
\begin{tabular}{cccccc}
\toprule
$\rho$ & E1 & E2 & E3 (makale) & E4 \\
\midrule
0   & 0.051 (\kod{rcond}~0.01 ile 3.87) & 0.051 & 0.017 & 0.017 \\
0.5 & 4.50 & 0.30 & 3.31 & 0.009 \\
1   & 4.88 & 0.115 & 4.39 & 0.012 \\
\bottomrule
\end{tabular}
\end{table}

\begin{table}[H]\centering\small
\caption{Sahte EDM $C$ [hedefin katı], 3 seed, 14~nm gürültü, 5~$\mu$m simetrik sürüklenme.
\kod{adim1\_edm.jsonl}~\BEL.}
\begin{tabular}{lccc}
\toprule
 & $\rho=0$ & $\rho=0.5$ & $\rho=1$ \\
\midrule
düzeltme yok & \multicolumn{3}{c}{5.0, 6.1, 9.5} \\
E1 & 0.003, 0.001, 0.001 & 13.7, 1.5, 0.95 & 6.9, 4.3, 0.36 \\
E2 & (=E1) & 0.20, 1.08, 0.26 & 0.004, 0.005, 0.002 \\
E3 (makale) & 0.001, 0.0002, 0.0003 & 0.27, 0.0001, 0.008 & 0.54, 3.1, 2.5 \\
E4 & (=E3) & 0.0002, 0.003, 0.005 & 0.0001, 0.0007, 0.005 \\
\bottomrule
\end{tabular}
\end{table}

\begin{figure}[H]\centering
\includegraphics[width=0.92\textwidth]{adim1_recon_k0_sym.pdf}
\caption{Kalan simetrik ve antisimetrik artık, BPM--kuadrupol korelasyonu $\rho$'ya karşı. Düz çizgi: her
yöntemin en iyi TSVD eşiği; noktalı: \kod{rcond}~$=0.01$.}
\end{figure}

\subsection{Model hatası ($\beta$-beat)}
Gerçek makineye \%0.2 ve \%1 rms gradyan hatası koyup (yaklaşık \%1 ve \%5 $\beta$-beat) kestiricileri nominal
modelle çalıştırdım~\BEL. 14~nm gürültü, $\rho=0$, 5~$\mu$m sürüklenmede iki demetin tek demete üstünlüğü
şöyle değişti: hatasız modelde 2.9 kat, \%0.2 hatada 1.2 kat, \%1 hatada 0.88 kat (yani tek demet daha
iyi). 1~$\mu$m gürültüde üstünlük her durumda yaklaşık 3 kat kaldı. \%0.2 hatayla sahte EDM ölçümü de
sonucu değiştirmedi: $\rho=0$'da E1 ve E3 ikisi de $C\le0.007$~\BEL.

\subsection{Yorum}
\begin{itemize}
\item Makaledeki yöntem (E3) yalnızca $\rho\approx0$'da çalışıyor. $\rho\ge0.5$'te simetrik artığın
büyük kısmı kalıyor; $\rho=1$'de 3 seed'in ikisinde $C>1$~\BEL.
\item $\rho=1$'de $\rho$'yu bilen tek demet (E2) hedefin 200 kat altına iniyor. Yani BPM kuadrupole
bağlıysa ``tek demet de görür'' öngörüsü doğru; ama standart yörünge düzeltmesi (E1) sahte EDM'yi
artırıyor~\BEL.
\item $\rho=0$'da bile ``tek demet çok zayıf'' iddiası büyük ölçüde sabit \kod{rcond}~$=0.01$
seçiminden geliyor: uygun eşikle tek demet simetrik sürüklenmeyi 5~$\mu$m'den 0.05~$\mu$m'ye indiriyor~\BEL.
\item İki demetin üstünlüğü, gürültünün model hatasından baskın olduğu rejimde var~\YOR.
\end{itemize}

%======================================================================
\section{Adım 2: Saptırıcı hataları (Hakem A, madde 1)}
%======================================================================

Hakem A, 48 saptırıcı için $10^{-4}$ bağıl alan hatası ve 0.1~mrad roll varsayıyor ve tek bir saptırıcının
yaklaşık 400~$\mu$m yörünge bozulması yaratacağını öngörüyor. Bu bölümde bu sayıyı doğruluyorum ve
saptırıcı sürüklenmesinin iki demetli izlemeyi nasıl etkilediğini inceliyorum.

\subsection{Neden önemli: elektrik kuvveti yönden bağımsızdır}
Manyetik kuvvet $q\,\vec v\times\vec B$ demet yönü çevrilince işaret değiştirir; elektrik kuvveti
$q\vec E$ değiştirmez. Kuadrupol kaçıklığı manyetik bir etkidir, saptırıcı hatası elektrik. Bu yüzden
saptırıcı hatasının yarattığı yörünge iki demette \emph{aynı işaretlidir} ve tam olarak toplam kanalına,
yani simetrik bileşeni okuduğumuz kanala düşer.

\subsection{Model}
Yeni motorda her saptırıcıya iki hata eklendi: bağıl alan hatası $\varepsilon$ ($E_0\to E_0(1+\varepsilon)$)
ve elektrik roll'u $\phi$. Roll, kapasitörün alan deseninin referans yörünge çizgisi etrafında
döndürülmesiyle uygulandı; öncü terimi düşey alan $E_z\approx E_0\phi$'dir.

\subsection{Hakem A'nın sayısı}
\begin{table}[H]\centering\small
\caption{Tek saptırıcı (indeks 10) hatasının kapalı yörüngeye etkisi. \kod{adim2\_kick.json}~\BEL.}
\begin{tabular}{lccc}
\toprule
hata & CCW rms / max [$\mu$m] & CW rms / max [$\mu$m] & corr(CW,CCW) \\
\midrule
$\varepsilon=10^{-4}$ (yatay) & 248 / 418 & 248 / 412 & $+0.95$ \\
elektrik roll 0.1~mrad (dikey) & 333 / 525 & 333 / 525 & $+0.95$ \\
eski manyetik \kod{dipole\_tilt} 0.1~mrad & 332 / 523 & 332 / 523 & $-0.95$ \\
\bottomrule
\end{tabular}
\end{table}
Hakemin $\sim$400~$\mu$m tahmini doğrulandı. Elektrik hatası iki demette aynı işaretli; eski manyetik
model işaret çeviriyor~\BEL.

\subsection{Fit matrisine saptırıcı sütunları}
Her düzlemde 48 saptırıcı sütunu eklenince bilinmeyen sayısı 96, denklem sayısı da 96 olur (2 demet
$\times$ 48 BPM). Rank tam kaldı (96/96), ama koşul sayısı x düzleminde 73'ten $5.6\times10^{3}$'e, y
düzleminde 118'den $7.9\times10^{3}$'e çıktı. Saptırıcı alt uzayı ile simetrik kuadrupol alt uzayı
arasında 4--6 doğrultuda asal açı kosinüsü 0.99'u aşıyor~\BEL (\kod{adim2\_fit.json}). Yani tam
dejenerelik yok, ama yakın dejenerelik var.

\subsection{Sürüklenme taraması}
\begin{figure}[H]\centering
\includegraphics[width=0.92\textwidth]{adim2_scan.pdf}
\caption{Saptırıcı sürüklenmesi altında simetrik kestirim hatası. Daire: makaledeki fit (yalnız kuadrupol
sütunları). Elmas: saptırıcı sütunları eklenmiş fit. 200 seed medyanı.}
\end{figure}
\begin{itemize}
\item Makaledeki fit, $10^{-7}$ alan sürüklenmesi ya da 0.1~$\mu$rad roll başına yaklaşık 1~$\mu$m simetrik
hata üretiyor~\BEL.
\item Saptırıcı sütunlu fit bu sızmayı gideriyor: kusursuz modelde hata 0.036~$\mu$m'de sabit; \%0.2
gradyan hatasıyla $\varepsilon=10^{-5}$'te 0.41~$\mu$m, $\phi=10~\mu$rad'da 0.92~$\mu$m~\BEL.
Bedeli, 1~$\mu$m gürültüde gürültünün yaklaşık 2 kat büyümesi (2.6 / 1.2~$\mu$m)~\BEL.
\end{itemize}

\subsection{Saptırıcı sürüklenmesinin kendi ürettiği sahte EDM}
\paragraph{Birinci mertebe roll etkisi.} Roll tek başına çok büyük bir sinyal üretiyor ve bu sinyal
genlikle \emph{doğrusal}:
\begin{table}[H]\centering\small
\caption{\kod{adim2\_roll\_C.json}~\BEL.}
\begin{tabular}{ccc}
\toprule
roll rms & $f_{\rm CW}=f_{\rm CCW}$ [hedef] & $C$ [hedef] \\
\midrule
$3\times10^{-9}$ & $-1.26\times10^{6}$ & 0.006 \\
$3\times10^{-8}$ & $-1.26\times10^{7}$ & 0.001 \\
$3\times10^{-7}$ & $-1.26\times10^{8}$ & 0.008 \\
\bottomrule
\end{tabular}
\end{table}
Mekanizmayı adım adım yazalım. (1) Sihirli momentumda elektrik alanı spini momentuma göre döndürmez; spin
düşey momentum dönmesini birebir izler. (2) Kapalı yörüngenin kapanması için düşey elektrik tekmelerini
kuadrupollerin manyetik tekmesi dengelemelidir: $\sum\theta_B=-\sum\theta_E$. (3) Manyetik alan ise spini
momentumdan $G\gamma$ kat fazla döndürür. Bu yüzden tur başına net spin dönmesi $G\gamma\sum\theta_B$ olur.

Sayısal kontrol: saptırıcı başına tekme $\theta_E=(L/R)\,\phi=0.131\,\phi$. 48 rastgele roll'un toplamı
$\sqrt{48}\times0.131\times3\times10^{-7}\approx2.7\times10^{-7}$~rad. $G\gamma\approx2.24$ ile tur
başına $\approx6\times10^{-7}$~rad, tur süresi $\approx3.7~\mu$s olduğundan $\approx0.16$~rad/s, yani
$\approx1.6\times10^{8}\times$hedef. Simülasyon $1.26\times10^{8}$ veriyor~\BEL. Bu etki CW$-$CCW farkında
$10^{-10}$ bağıl hassasiyetle sadeleşiyor ($C$ sayısal tabanda); ama gerçek makinede bu sadeleşme CW/CCW
simetrisine bağlıdır~\YOR.

\paragraph{$\varepsilon\times\phi$ çarpım etkisi.} $\varepsilon$ ve $\phi$ birlikte $3\times10^{-7}$
olduğunda, kuadrupoller kusursuzken bile $C=1.17$, 0.79, 2.09 kalıyor~\BEL. Bu, saptırıcıların ürettiği
yatay ve düşey kapalı yörüngelerin çarpımından ($x_{\rm CO}\,y_{\rm CO}$) gelen geometrik fazdır;
kuadrupol kaçıklığındaki mekanizmanın aynısı~\YOR. Yalnız kuadrupolleri düzeltmek bunu gidermiyor
($C=1.20$, 0.81, 2.51)~\BEL. Çözüm Bölüm~\ref{sec:T1}'de.

%======================================================================
\section{Adım 3: Düzelticiler}
%======================================================================

Hakem A elektrostatik ve manyetik düzelticilerin farkının açıklanmasını istiyor. Bu bölümde aynı iki
demetli düzeltmeyi dört düzeltici türüyle koştum.

\subsection{Model}
\begin{center}\small
\begin{tabular}{llll}
\toprule
kısaltma & tür & konum & kanal \\
\midrule
Mq & manyetik (makalenin örtük modeli) & kuadrupol içinde & fark \\
Md & manyetik dipol & saptırıcıda & fark (corr $-0.95$) \\
Ed & elektrostatik ($\varepsilon$, $\phi$ trimi) & saptırıcıda & toplam (corr $+0.95$) \\
M+Ed & ikisi birlikte & saptırıcıda & ikisi \\
\bottomrule
\end{tabular}
\end{center}
İki hedef karşılaştırıldı: ``sıfır'' (BPM okumasını sıfırla) ve ``merkez'' (demeti kestirilen kuadrupol
merkezlerinden geçir). Bu ayrım önemlidir: laboratuvara sabit BPM'de yörüngeyi sıfırlamak, demetin
sürüklenmiş kuadrupollerin merkezinden geçtiği anlamına gelmez.

\subsection{Sonuçlar}
\begin{table}[H]\centering\small
\caption{Senaryo Adım 1 ile aynı. 3 seed. \kod{adim3\_edm.jsonl}~\BEL.}
\begin{tabular}{lccc}
\toprule
senaryo & demet $-$ kuad. merkezi [$\mu$m] & $C$ & çift kısım \\
\midrule
düzeltme yok & 7.5--10.6 & 5.0, 6.1, 9.5 & 0.8--10 \\
\textbf{Mq} & 0.02--0.03 & \textbf{0.0002, 0.0005, 0.0013} & $\le$0.004 \\
Md, sıfır & 5.0 & 0.04, 1.01, 0.76 & $\sim$200 \\
Md, merkez & 3.6--4.1 & 55, 6.1, 12 & 15--700 \\
Ed, sıfır & 4.8 & 14, 4.4, 6.7 & $2\times10^{7}$--$1.5\times10^{8}$ \\
Ed, merkez & 0.07--0.12 & 2.9, 1.5, 2.8 & $10^{6}$--$9\times10^{6}$ \\
M+Ed, sıfır & 5.0 & 2.2, 16, 18 & $\sim10^{8}$ \\
M+Ed, merkez & 0.03--0.21 & 5.5, 4.4, 1.1 & $\sim10^{6}$--$10^{7}$ \\
\bottomrule
\end{tabular}
\end{table}

\begin{figure}[H]\centering
\includegraphics[width=0.92\textwidth]{adim3_edm.pdf}
\caption{Düzeltici türlerine göre $C$ ve çift kısım.}
\end{figure}

\subsection{Yorum}
\begin{itemize}
\item Yalnız kuadrupol içi düzeltici (Mq) çalışıyor, çünkü kaçıklığın dipol alanını aynı yerde siler:
$g(y-d)+\Delta B=g\,y$, eğer $\Delta B=g\,d$~\YOR.
\item Simetrik sürüklenme, demetin iki yönde \emph{aynı işaretle} kaydırılmasını gerektirir. Manyetik
düzeltici fark kanalına etki ettiği için bunu yapamaz; elektrostatik düzeltici yapabilir (demet 0.1~$\mu$m
yakından geçer), ama birinci mertebe çift sinyal ve $C\approx1.5$--3 artık bırakır~\BEL.
\item Elektrostatik düzelticinin spini doğrudan döndürmediği doğrudur; ama yörüngenin kapanması için
kuadrupollerin manyetik tekmesini zorunlu kılarak sahte EDM'yi o yoldan üretir~\YOR.
\end{itemize}

%======================================================================
\section{Adım 4: Korelasyonlu hizalama hataları (Hakem A, madde 2)}
%======================================================================

Hakem A, iki kuadrupol arasındaki göreli hatanın aradaki mesafenin kareköküyle büyüdüğü bir model
öneriyor. Bu bölümde böyle bir modelde simetrik bileşenin baskınlaşıp baskınlaşmadığını ve Şekil 4'ün
nasıl değiştiğini inceliyorum.

\subsection{Model}
Kuadrupoller eşit aralıklıdır: $\Delta s=16.66$~m, çevre 799.7~m.
\begin{itemize}
\item \textbf{Rastgele yürüyüş:} $\langle(y_i-y_j)^2\rangle=a^2|s_i-s_j|$. Halka kapalı olduğu için
periyodik Brownian köprüsü kullanıldı.
\item \textbf{ATL:} Shiltsev'in yer hareketi yasası $\langle\Delta Y^2\rangle=A\,T\,L$, sabit $T$ için
uzayda rastgele yürüyüşle aynı yapıyı verir; farkı genliğin zamanla $\sqrt{T}$ büyümesidir. Adım 6'da
kullanıldı.
\item \textbf{Karma:} rastgele yürüyüş + yerel bağımsız hata.
\end{itemize}
Rijit öteleme ve eğim demet tarafından izlendiği için zararsızdır; dikeyde $m=0,1$, yatayda $m=1$
harmonikleri çıkarıldı ve desenler 10~$\mu$m rms'e normalize edildi.

\subsection{Simetrik/antisimetrik ayrışım}
\begin{table}[H]\centering\small
\caption{2000 seed, iki düzlem. \kod{adim4\_stats.json}~\BEL.}
\footnotesize
\begin{tabular}{lccc}
\toprule
model & sim/antisim rms oranı [16--84\%] & simetrik varyans kesri & komşu farkı \\
\midrule
korelasyonsuz & 0.94 [0.76, 1.15] & 47\% & 14.6~$\mu$m \\
rastgele yürüyüş (ham) & 3.6 [2.6, 5.0] & 93\% & 5.4~$\mu$m \\
rastgele yürüyüş (düzlem çıkarılmış) & 2.3 [1.8, 3.0] & 84\% & 8.0~$\mu$m \\
karma (düzlem çıkarılmış) & 1.7 [1.4, 2.2] & 75\% & 10.1~$\mu$m \\
\bottomrule
\end{tabular}
\end{table}
Beklendiği gibi, komşu kuadrupoller birlikte hareket ettiği için simetrik bileşen baskınlaşıyor~\BEL.

\subsection{Şekil 4'ün tekrarı}
\begin{figure}[H]\centering
\includegraphics[width=0.92\textwidth]{adim4.pdf}
\caption{(a) Harmonik içerik. (b) Sahte EDM: ham, tek demet ve iki demetli düzeltme (laboratuvara sabit
BPM, 14~nm, kuadrupol içi düzeltici), 3 seed.}
\end{figure}
Ham sahte EDM: korelasyonsuz 294, 56, 124; rastgele yürüyüş 364, 399, 203; karma 523, 118, 2265.
Düzeltme sonrası her modelde $C\le0.005$~\BEL. Model tam bilinirken kestirici neredeyse tam ters
aldığından kalan artık yalnızca gürültüdür ve korelasyon modelinden bağımsızdır~\YOR.

\subsection{Yer hareketi BPM'leri de taşırsa}
Gerçekte zemin hem kuadrupolleri hem BPM'leri taşır. Kuadrupol $=$ zemin $+$ yerel, BPM $=$ zemin $+$ yerel
modeliyle (toplam 5~$\mu$m, 14~nm):
\begin{table}[H]\centering\small
\caption{Kalan simetrik rms [$\mu$m]. \kod{adim4\_ground\_bpm.json}~\BEL.}
\begin{tabular}{ccccc}
\toprule
yerel/zemin & E1 & E3 (makale) & E4 (BPM'in zemini izlediğini bilen) \\
\midrule
0   & 4.70 & 3.31 & 0.010 \\
0.1 & 4.69 & 3.31 & 0.37 \\
0.3 & 4.58 & 3.28 & 1.03 \\
1   & 4.09 & 3.20 & 2.52 \\
\bottomrule
\end{tabular}
\end{table}
Uzun dalgalı yer hareketinde BPM'ler kuadrupollerle birlikte hareket eder ($\rho\approx1$). Bu durumda
makaledeki model simetrik sürüklenmenin yalnızca yaklaşık \%30'unu yakalar~\BEL. Hakem A'nın endişesinin
asıl karşılığı budur~\YOR.

%======================================================================
\section{Tartışma testleri: yöntemi kurtaran iki dokunuş}
%======================================================================

Adım 1--4 sonunda iki sorun kaldı: saptırıcı sürüklenmesi ve BPM hareketi. İkisini aynı ilkeyle ele aldım:
\emph{hatayı kendi kaynağında ve aynı tür alanla gör ya da düzelt}.

\subsection{T1: Saptırıcıyı kaynağında trimlemek}\label{sec:T1}
Saptırıcı sütunlu fitle her saptırıcının $\Delta\varepsilon$ ve $\Delta\phi$ değerini kestirip, alan
hatasını yüksek gerilim trimiyle, roll'u ise üst ve alt plakalar arasına uygulanan küçük bir gerilim farkıyla
(düşey elektrik trimi) geri aldım.

\begin{table}[H]\centering\small
\caption{Adım 2 senaryosu, 3 seed. \kod{t1\_edm.jsonl}~\BEL.}
\begin{tabular}{lcc}
\toprule
durum & $C$ & çift kısım \\
\midrule
yalnız kuadrupol düzeltmesi & 1.20, 0.81, 2.51 & $\sim10^{8}$ \\
\textbf{kuadrupol + saptırıcı trimi} & \textbf{0.0010, 0.0005, 0.0017} & $6\times10^{3}$--$2.5\times10^{6}$ \\
kontrol: saptırıcı sürüklenmesi tam biliniyor & 0.0002, 0.0050, 0.0053 & $\sim10^{-3}$ \\
\%0.2 gradyan hatası, trim & 0.010, 0.011, 0.006 & $\sim2\times10^{6}$ \\
\%0.2 gradyan hatası, yalnız kuadrupol & 1.28, 0.74, 2.80 & $\sim10^{8}$ \\
\bottomrule
\end{tabular}
\end{table}

Neden çalıştığını açıklayayım. Tek tek saptırıcıların kestirimi kötüdür: trimden sonra hata başlangıcın
\%25--80'i kadar kalır. Ama kötü kestirilen doğrultular, yörüngeye zaten küçük katkı veren doğrultulardır.
Saptırıcı kaynaklı \emph{yörünge} \%97--99 azalır~\BEL. $\varepsilon\times\phi$ etkisi iki yörüngenin
çarpımıyla orantılı olduğu için bu yeterlidir~\YOR. Roll trimi, öncü terim olan düşey alanı ekleyerek
modellendi; alan deseninin dönmesinden gelen ikinci mertebe fark ihmal edildi~\HIP.

\subsection{T2: BPM'i kuadrupole rijit bağlamak}
BPM kuadrupole rijit bağlıysa okuma $r=x-d$ olur, yani demetin kuadrupol merkezine göre konumu. Bu,
sahte EDM'yi belirleyen büyüklüğün ta kendisidir. Geriye kalan hata kaynağı, BPM elektriksel merkezi ile
kuadrupol manyetik merkezi arasındaki farkın BBA'dan sonraki değişimi $\delta$'dır (mekanik kayma ve BPM
elektroniğinin ofset kayması toplamı).

\begin{figure}[H]\centering
\includegraphics[width=0.92\textwidth]{t2.pdf}
\caption{Rijit bağlı BPM: kalan simetrik artık, göreli kayma $\sigma_\delta$'ya karşı.}
\end{figure}

\begin{table}[H]\centering\small
\caption{Sahte EDM $C$, 3 seed, 5~$\mu$m sürüklenme, 14~nm. \kod{t2\_edm.jsonl}~\BEL.}
\begin{tabular}{ccc}
\toprule
$\sigma_\delta$ & iki demet (E4) & tek demet (E2) \\
\midrule
0.1~$\mu$m & 0.001, 0.003, 0.005 & 0.04, 0.10, 0.004 \\
0.3~$\mu$m & 0.012, 0.036, 0.026 & 0.63, 0.28, 0.40 \\
1~$\mu$m & 0.17, 0.35, 0.27 & 8.8, 10.1, 3.2 \\
\bottomrule
\end{tabular}
\end{table}
İki demetin hatası yaklaşık $1.1\,\sigma_\delta$ ve sürüklenme genliğinden neredeyse bağımsız. Bu mimaride
iki demetin üstünlüğü 25--50 kat~\BEL.

\textbf{Literatür.} BPM'i kuadrupole bağlamak lineer hızlandırıcılarda yerleşik bir uygulamadır: DESY
TTF2/European XFEL'de BPM'ler kuadrupollerin içindedir ve gerili telle manyetik eksene kalibre edilir;
CLIC/PACMAN'da BPM ve kuadrupol ortak desteğe mikrometre hassasiyetle monte edilir; LHC ark BPM'leri
kuadrupol kriyostatlarındadır~\BEL. Işık kaynakları ise senkrotron ışımasının ısıttığı girder yüzünden
BPM'leri tersine ayırır (SLS, NSRRC)~\BEL. Proton EDM halkasında senkrotron ışıması ihmal edilebilir, bu
yüzden lineer hızlandırıcı gerekçesi geçerlidir~\YOR. Ancak literatür µm düzeyinde \emph{ön hizalama}
gösterir; bizim istediğimiz zamanla kararlılıktır. Bu yüzden $\sigma_\delta\lesssim0.3$--1~$\mu$m bir
tasarım gereksinimi olarak yazılmalıdır~\YOR.

%======================================================================
\section{Adım 5: Harmonik başına yörünge kazancı ve spin duyarlılığı}
%======================================================================

Hakem A, uzun dalgalı kaymada demet kuadrupolleri izlediği için yörüngenin bastırıldığını, ama spinin aynı
bastırmayı görüp görmediğinin belli olmadığını söylüyor. Bu bölümde tek harmonikli desenlerle bunu
ölçüyorum.

\subsection{Model}
Hücre harmoniği $m=0\ldots12$ için simetrik desen $v_{2k}=v_{2k+1}=\cos(2\pi mk/24)$, antisimetrik desen
$v_{2k}=-v_{2k+1}=\cos(2\pi mk/24)$. Yörünge kazancı $\|Rv\|/\|v\|$. Sahte EDM bilineer olduğu için $dx$ ve
$dy$ aynı harmonikle ve aynı türle eşleştirildi; iki eşleşme kullanıldı: \emph{aynı faz} ($dx=dy=\cos$) ve
\emph{dik faz} ($dx=\cos$, $dy=\sin$; $m=0,12$'de tanımsız). Genlik 5~$\mu$m; sonuç $C/\mu$m$^2$.

\begin{figure}[H]\centering
\includegraphics[width=0.92\textwidth]{adim5.pdf}
\caption{Yörünge kazancı (sol eksen) ve birim genlik karesi başına sahte EDM (sağ eksen).}
\end{figure}

\begin{table}[H]\centering\small
\caption{$C$ per $\mu$m$^2$ ve dikey yörünge kazancı. \kod{adim5\_edm.jsonl}, \kod{adim5\_gain.json}~\BEL.}
\begin{tabular}{cccccc}
\toprule
$m$ & sim, aynı & sim, dik & antisim, aynı & antisim, dik & kazanç sim/antisim \\
\midrule
0 & $9\times10^{-5}$ & -- & 0.83 & -- & 1.01 / 6.9 \\
1 & 0.12 & 0.84 & 16 & 0.20 & 1.4 / 8.5 \\
2 & 3.2 & 11 & 106 & 0.99 & 5.8 / 29 \\
3 & 4.0 & 9.0 & 59 & 6.4 & 2.4 / 9.6 \\
6 & 0.16 & 0.15 & 0.52 & 0.040 & 0.55 / 1.16 \\
10 & 0.070 & 0.019 & 0.082 & 0.0024 & 0.35 / 0.40 \\
12 & 0.0324 & -- & 0.0324 & -- & 0.34 / 0.34 \\
\bottomrule
\end{tabular}
\end{table}

\subsection{Yorum}
\begin{itemize}
\item $m\approx12$'de simetrik ve antisimetrik desen ayırt edilemez: kazançlar ve sahte EDM birebir aynı.
Düşük $m$'de simetrik kazanç antisimetrikten yalnızca 4--7 kat küçük~\BEL. Makaledeki ``$k=24$ bileşeni
yaklaşık 100 kat bastırılır'' ifadesi ölçülen kazançlarla uyuşmuyor~\YOR.
\item Spin, $m=0$'da çok güçlü bastırılır ($9\times10^{-5}$): demet kuadrupolleri izleyip merkezlerinden
geçer. $m=1$'de bastırma kısmi, en büyük duyarlılık betatron tune'una yakın $m=2$--3'te~\BEL.
Korelasyonlu hata modelleri gücünü tam bu harmoniklerde toplar (Adım 4)~\YOR.
\item Yörünge kazancının karesine bölünmüş en büyük sahte EDM, simetrik bileşen için antisimetrikten
2--3 kat büyüktür ($m=2$'de 0.33'e karşı 0.13)~\BEL. Yani simetrik bileşen yörüngedeki görünürlüğüne göre
spine daha duyarlıdır, ama fark ılımlıdır.
\end{itemize}

%======================================================================
\section{Adım 6: Zaman senaryosu (Hakem B)}
%======================================================================

Hakem B, yöntemin işleyişinin ve performansının titiz bir biçimde gösterilmesini istiyor. Bu bölümde bir
haftalık gerçekçi bir işletim senaryosu kuruyorum.

\subsection{Senaryo}
\begin{itemize}
\item 1 hafta, 1 saatlik adımlar. Sürüklenme: ATL yer hareketi, $A=10^{-5}~\mu$m$^2$/(s$\cdot$m); LEP
($1.1\times10^{-5}$) ve TRISTAN ($2.7\times10^{-5}$) ölçümlerine göre seçildi~\BEL.
\item $t=0$'da BBA artığı 0.3~$\mu$m rms~\HIP. BPM--kuadrupol kayması 0.2~$\mu$m$/\sqrt{\text{gün}}$~\HIP;
beklenen kayma 0.3~$\mu$m'ye ulaşınca BBA tetiklenir (bu hızla yaklaşık her 54 saatte).
\item Gürültü: 10~$\mu$m/tur, 60~s ortalama, 10~nm elektronik tabanı~\HIP. Ölçüm başına
$\sigma_n=\sqrt{\sigma_{\rm tur}^2/(f_{\rm dev}T_{\rm ort})+\sigma_{\rm taban}^2}\approx10$~nm. Burada
$f_{\rm dev}\approx270$~kHz tur frekansıdır.
\item Senaryolar: W1 iki demet (rijit BPM, E4), W1 tek demet (E2), W0 makale modeli (BPM'ler ayrı ayakta ama
aynı zeminde, E3), yalnız periyodik BBA, hiçbir şey.
\end{itemize}

\begin{figure}[H]\centering
\includegraphics[width=\textwidth]{adim6.pdf}
\caption{(a) Simetrik artık, doğrusal, 20 seed. (b) Sahte EDM, C++, 3 seed. (c) İki demetli tespit
sınırı.}
\end{figure}

\begin{table}[H]\centering\small
\caption{Sahte EDM $C$, 3 seed (küçükten büyüğe). \kod{adim6\_edm.jsonl}~\BEL. Hız için kısaltılmış
estimatör ayarı kullanıldı.}
\footnotesize
\begin{tabular}{lccccc}
\toprule
senaryo & 6 sa & 24 sa & 48 sa & 96 sa & 168 sa \\
\midrule
$t=0$ (BBA artığı) & \multicolumn{5}{l}{0.07 / 0.73 / 1.35} \\
\textbf{W1, iki demet} & 0.05/0.84/1.3 & 0.16/0.37/1.1 & 0.10/0.31/1.4 & 0.22/0.69/0.95 & 0.12/0.31/0.34 \\
W1, tek demet & 0.16/1.2/2.4 & 0.72/0.98/6.2 & 0.39/2.6/7.6 & 2.8/4.7/12 & 0.33/0.95/1.0 \\
W0, makale & 0.89/2.2/3.9 & 1.6/2.7/5.5 & 0.8/7.8/13 & 7.5/15/21 & 0.6/1.9/4.2 \\
yalnız BBA & 4.3/9.5/43 & 2.0/5.7/34 & 28/32/112 & 13/122/129 & 3.5/7.8/15 \\
hiçbir şey & 4.3/9.5/43 & 2.0/5.7/34 & 28/32/112 & 97/501/673 & 232/427/718 \\
\bottomrule
\end{tabular}
\end{table}

\subsection{En küçük tespit edilebilir sürüklenme}
Simetrik kestirimin gürültü kazancı iki demette $k=0.75$, tek demette 2.35, makale modelinde 1.23~\BEL.
3$\sigma$ tespit sınırı $3k\sigma_n$'dir. Örnek hesap: 10~$\mu$m/tur ve 10~nm taban, $T_{\rm ort}=1$~s
için $\sigma_n=\sqrt{(10~\mu\text{m})^2/(270\,000\times1)+(10~\text{nm})^2}=\sqrt{(19.2)^2+10^2}$~nm
$\approx21.7$~nm; sınır $3\times0.75\times21.7\approx49$~nm $\approx0.05~\mu$m. Ortalama süresi uzayınca
sınır elektronik tabanına iner: 10~nm tabanla $\approx0.02~\mu$m, 100~nm tabanla $\approx0.23~\mu$m~\BEL.

\subsection{Yorum}
\begin{itemize}
\item İki demetli izleme bir hafta boyunca sahte EDM'yi BBA artığının belirlediği seviyede tutuyor
(0.05--1.4); yalnız periyodik BBA iki BBA arasında 10--100 kat yükseliyor~\BEL.
\item İzleyici BBA'nın bıraktığı durumu korur, daha iyisini yapamaz. Hedefin güvenli biçimde altında kalmak
için BBA artığı $\lesssim0.15$--0.2~$\mu$m olmalı~\HIP (sahte EDM'nin $\sigma^2$ ölçeklemesinden).
\item BPM'ler zemini izlerken makaledeki model hedefin altında kalamıyor~\BEL.
\end{itemize}

%======================================================================
\section{Makale için çıkarımlar}
%======================================================================

\subsection{Yöntemin geçerlilik alanı}
\YOR~Makalenin bugünkü iddiası kabaca ``iki demetli toplam, kör simetrik bileşeni görür ve sürüklenme
izleyicisi olarak kullanılabilir'' şeklindedir. Kontroller bu iddianın aşağıdaki daha dar ama savunulabilir
biçimini destekliyor:

\begin{quote}
BPM'ler kuadrupollere rijit bağlı ve kestirici bunu bildiğinde, düzelticiler kuadrupol üstü manyetik trim
sargıları olduğunda ve saptırıcılar kendi elektrik trimleriyle düzeltildiğinde, iki demetli izleme tek
demetten 25--50 kat iyi çalışır ve BBA kalitesini günler boyunca korur.
\end{quote}

\subsection{Tasarım gereksinimleri}
\begin{center}\small
\begin{tabularx}{\textwidth}{lX}
\toprule
gereksinim & dayanak \\
\midrule
BPM--kuadrupol göreli kararlılığı $\sigma_\delta\lesssim0.3$--1~$\mu$m (BBA aralığı boyunca) & T2 \\
kuadrupol üstü manyetik dipol trim sargıları & Adım 3 \\
saptırıcı başına yüksek gerilim ve düşey elektrik trimi; saptırıcı sütunlu fit & Adım 2, T1 \\
yalnız kuadrupol fiti kullanılırsa saptırıcı kararlılığı $\lesssim5\times10^{-8}$ ve 0.05~$\mu$rad & Adım 2 (\HIP\ eşik) \\
BBA artığı $\lesssim0.15$--0.2~$\mu$m & Adım 6 (\HIP) \\
BPM elektronik tabanı $\lesssim10$--100~nm & Adım 6 \\
çift kısmın sadeleşmesi için CW/CCW simetrisi & Adım 2, T1 \\
\bottomrule
\end{tabularx}
\end{center}

\subsection{Metinde düzeltilmesi gerekenler}
\begin{enumerate}
\item BPM modeli açıkça yazılmalı; Sec.~V.B'deki ``ortak hareket'' argümanı ters çevrilmeli: ortak hareket
yöntemin aleyhine değil, kestirici bunu bildiğinde lehine.
\item Düzeltici türü (kuadrupol içi manyetik) ve fitin düzenlileştirmesi (TSVD, \kod{rcond}) belirtilmeli.
\item Şekil 4 elmasının koşulları (14~nm, sıfır BPM ofseti, 15 seed) yazılmalı.
\item Sec.~III'teki 100 kat bastırma iddiası ve Şekil 1(a)'daki analitik $G_k$ formülü, ölçülen harmonik
tabloyla değiştirilmeli.
\item ``Tek demet kördür'' yerine nicel bir karşılaştırma verilmeli: laboratuvara sabit BPM'de yaklaşık
3 kat (model hatasıyla kaybolur), rijit BPM'de 25--50 kat.
\item Hakem B için Adım 6'nın zaman senaryosu ve tespit sınırı eklenmeli.
\end{enumerate}

%======================================================================
\section{Sınırlamalar ve açık sorular}
%======================================================================
\begin{itemize}
\item C++ sahte EDM sonuçlarının çoğu 3 seed'e dayanır; seed'ler arası dağılım bazen 10 kattan fazladır.
Makaleye girecek sayılar için daha fazla seed gerekir.
\item Statik başlangıç durumu simüle edilmedi: $10^{-4}$/0.1~mrad statik saptırıcı hataları milimetre
düzeyinde yörünge yaratır; bunu kalibrasyonda düzeltmenin kendi sahte EDM tabanı bilinmiyor~\HIP.
\item Adım 6'da saptırıcı sürüklenmesi ve T1 trimi aynı senaryoda birleştirilmedi.
\item Adım 6 parametreleri (BBA artığı, $\delta$ kayma hızı, elektronik taban) tahmindir.
\item Elektrostatik düzelticilerde kalan $C\approx1.5$--3'ün mekanizması doğrulanmadı.
\item Adım 6'da kısaltılmış estimatör ayarı kullanıldı (daha önce $p=2.00$ ile doğrulanmış).
\end{itemize}

%======================================================================
\section{Çalışma sırasında düzeltilen hatalar}
%======================================================================
\begin{enumerate}
\item \textbf{Adım 1, ilk sürüm:} $\rho=0$'da bile bağımsız BPM hareketi tam genlikle eklenmişti; böylece
``laboratuvara sabit'' durum yanlışlıkla ``BPM'ler de 5~$\mu$m sürükleniyor'' durumuna dönmüştü. İki aile
ayrıldı (u $=\rho\Delta m$ ve eşit rms'li kötümser aile).
\item \textbf{Adım 2, tekme testi:} Sıfır testinde eski motoru yükleyen işçi süreci bu motoru bellekte
tuttuğu için bir CCW roll koşumu yanlışlıkla sıfır çıkmıştı. Motor her görevde açıkça seçilecek biçimde
düzeltildi.
\item \textbf{Adım 5:} $m=0$ simetrik desen sabit olduğundan standart sapmayla normalizasyon NaN verdi;
rms ile normalize edildi. Ayrıca bir ara sohbette simetrik bileşenin spin/yörünge oranı ``yaklaşık 50 kat''
diye yanlış yazılmıştı (antisimetrik için yanlış eşleşmeyle kıyaslanmıştı). Doğrusu 2--3 kattır.
\item \textbf{T2:} Hesap yanlışlıkla iki kez başlatılmıştı; kopya durduruldu, yinelenen satırlar
temizlendi.
\end{enumerate}

%======================================================================
\appendix
\section{Yeniden üretim}
%======================================================================
Bütün betikler ve sonuçlar \kod{claude/brave-allen-92wn61} dalında, \kod{revizyon/} klasöründedir. Mevcut
kod değiştirilmedi.

\begin{center}\small
\begin{tabularx}{\textwidth}{lX}
\toprule
betik & işlev \\
\midrule
\kod{integrator\_rev.cpp/.py} & saptırıcı $\varepsilon$, $\phi$ ve saptırıcı konumunda manyetik düzelticiler \\
\kod{adim1\_bpm\_model.py} & BPM modeli, $\rho$ taraması, $\beta$-beat (\kod{ADIM1\_SIG\_G}) \\
\kod{build\_rtrue.py} & gradyan hatalı gerçek makine matrisleri \\
\kod{adim2\_deflector.py} & tekme testi, saptırıcı matrisleri, rank/tarama, sahte EDM \\
\kod{adim2\_defl\_diag.py}, \kod{adim2\_roll\_C.py} & roll teşhisi ve ölçekleme \\
\kod{adim3\_correctors.py} & düzeltici türleri \\
\kod{adim4\_correlated.py}, \kod{adim4\_ground\_bpm.py} & korelasyonlu hata ve zemin--BPM \\
\kod{t1\_defl\_trim.py}, \kod{t2\_bpm\_rigid.py} & tartışma testleri \\
\kod{adim5\_harmonics.py} & harmonik kazanç ve sahte EDM \\
\kod{adim6\_time.py} & zaman senaryosu ve tespit sınırı \\
\bottomrule
\end{tabularx}
\end{center}

\end{document}
